\documentclass[11pt]{article}
\usepackage[margin=1in]{geometry}
\usepackage{amsmath,amssymb,amsthm,mathtools}
\usepackage{enumitem}
\usepackage{booktabs}
\usepackage{microtype}
\usepackage{xcolor}
\usepackage{hyperref}
\usepackage{cleveref}
\hypersetup{colorlinks=true, linkcolor=blue, citecolor=blue, urlcolor=blue}

\newtheorem{definition}{Definition}

\title{EvidenceVault: zk-RTP: Evidence Packages and Support Manifests}
\author{}
\date{March 28, 2026}

\begin{document}
\maketitle

\begin{abstract}
Portable zk-RTP evidence in this family is carried by bundle-scope evidence packages and support manifests on one bundle-policy scope.
Evidence packages carry single-object bundle claims, while support manifests carry the tagged multi-object support family on that same scope.
\end{abstract}

\section{Evidence packages and support manifests}

\begin{definition}[Bundle-scope evidence package]
A bundle-scope evidence package is the family's typed local evidence carrier on that fixed bundle-policy scope and is the object
\[
E=(\Phi,q,p,\Sigma,\sigma),
\qquad
\Phi=(d_B,d_v,h_P,d_\nu,c),
\]
where $\Phi$ is the carried bundle-policy scope tuple, with $d_B$ the carried canonical digest of the family's local bundle $B$, $d_v$ the verifier / checker identity digest, $h_P$ the governing policy digest, $d_\nu$ any resolver or profile digest needed to decode local objects, and $c$ the carried checkpoint or time-anchor context; $q$ is the carried direct evidence kind and $p$ is the corresponding payload, $\Sigma$ is the carried bridge replay slice, and $\sigma$ authenticates the canonical encoding of the carried fields when the kind is signable.
The direct kind/payload fields then have one of the following local forms.
\begin{enumerate}[label=\textnormal{(\alph*)},leftmargin=2em]
\item \textbf{Acceptance case.} $q=\mathsf{accept}$ and $p$ records the local acceptance verdict together with any realized menu or branch choices (or invariance evidence for them).
\item \textbf{Claim case.} $q=\mathsf{claim}$ and
\[
p=(h_{\mathrm{schema}},\mathsf{claims}),
\]
where $h_{\mathrm{schema}}$ pins the claim schema and version and $\mathsf{claims}$ is the canonical claim family being carried.
\item \textbf{Freshness case.} $q=\mathsf{fresh}$ and $p=(d_e,t_0,\mathcal{R},\mathcal{W})$ binds one evidence digest $d_e$ to an authenticated commitment time $t_0$, an admissible source class $\mathcal{R}$, and the carried freshness window $\mathcal{W}$, asserting that the decision time satisfies the policy's admissible temporal predicate determined by $t_0$ and $\mathcal{W}$.
\item \textbf{Governed-action cases.} If
\[
q\in\{\mathsf{strengthen},\mathsf{waiver},\mathsf{override},\mathsf{renewal},\textsf{emergency-accept},\textsf{emergency-reject},\mathsf{revocation}\},
\]
then $p$ has the form
\[
p=\Gamma,
\]
whose local $\delta$-reading here is the canonical policy-diff or other action object relevant to that governed case.
\end{enumerate}
\end{definition}

\begin{definition}[Support manifest]
A support manifest is the family's support carrier on that fixed bundle-policy scope and is the tuple
\[
\mathsf{SM}=(\Phi,d_{\top},\mathcal{S}),
\]
where $\Phi$ is the carried shared bundle-policy scope tuple and $d_{\top}$ is the digest of the top-level decision or claim object the manifest is meant to support.
It carries one tagged family of support entries $s=(q,r,d,\varsigma)$, where $d$ is the digest of a member bundle-scope evidence package. If $q=\mathsf{portfolio}$, then $r$ is any carried support role and $\varsigma$ carries whatever dependency class, provenance label, or threshold-counting annotation the manifest declares for that member. If $q=\mathsf{dual}$, then $r\in\{\textsf{core},\textsf{boundary}\}$ and $\varsigma$ records that those designated roles are jointly required complementary obligations for the same supported decision or claim. Any realized portfolio member or dual witness therefore appears directly as one such entry in $\mathcal{S}$ rather than through a second family-local support-case surface.
\end{definition}

\end{document}
